\section{Discussion and Limitations}
\label{sec:discussion}

\paragraph{What IsoColBERT does and does not do.} The retrieval gains concentrate where two conditions intersect: weak pretrained signal and retrieval headroom. High-resource pairs sit at the P@1 ceiling for both methods, and the small ($-0.05$ pp on EN-ES, not significant) shift on the training pair shows that the regularizer does not damage in-distribution accuracy. The regularizer reshapes geometry in proportion to its weight $\lambda$, but the retrieval payoff saturates once anisotropy stops being the bottleneck. Beyond that point, additional regularization continues to spread the manifold without further accuracy gains.

\paragraph{Asymmetric collapse under bitext contrastive training.} An incidental finding that we have not seen reported elsewhere: under symmetric bitext InfoNCE, the target-side language collapses harder than the source-side language. Spanish tokens have intra-cosine $0.29$ under vanilla training while English tokens have $0.20$; effective ranks are $101.7$ and $105.1$ respectively. We attribute this to the geometry of pretraining: XLM-R is exposed to far more English than Spanish text, so the contrastive objective has more work to do on the Spanish side and the corresponding representations bear most of the gradient pressure. IsoColBERT closes the gap. The geometry of training-pair language asymmetry suggests a diagnostic for predicting which side of a contrastive setup will benefit most from explicit isotropy pressure.

\paragraph{Relation to I-STAR.} \citet{rudman2024istar} report that for classification fine-tuning, decreasing isotropy can improve accuracy. Our result on retrieval is in the opposite direction. The settings differ in two ways that we believe explain the disagreement. First, classification benefits from local clustering: examples of the same class should have nearby representations, which is locally anisotropic. Retrieval benefits from angular separation across the full manifold so that the argmax in MaxSim has high cosine margin. Second, classification operates on sentence-pooled representations, while late-interaction retrieval operates on individual token vectors; the geometric quantity that matters is different. We take this as evidence that the productive direction for isotropy regularization is task-dependent and that retrieval is a clean case where higher isotropy across the embedding manifold is the right objective.

\paragraph{Limitations.} We train on a single language pair (EN-ES) and evaluate transfer to six others. We do not test whether multi-pair training, mixed objectives, or curriculum strategies would further improve the low-resource gains. The backbone is XLM-R base; we do not test whether the effect transfers to larger multilingual encoders or to retrieval-specialized backbones like mE5. The benchmark family is FLORES+; we do not yet test on MIRACL \citep{zhang2023miracl} or mMARCO \citep{bonifacio2021mmarco}, both of which are harder and more realistic retrieval benchmarks. Three seeds is the practical minimum for reporting variance, and we note the per-seed values throughout to allow audit. The $\lambda$ grid is coarse ($\{0, 0.1, 0.5\}$); finer-grained search around $\lambda=0.5$ may yield additional retrieval payoff. Finally, we present mechanism ablations that rule out two alternative regularizer placements, but other placements (for example, regularizing only document tokens, or restricting to a percentile of tokens) remain unexplored.
